\clearpage\label{answers}
{\sans\bfseries\LARGE\color{ink} Answers}\par
\footnotesize\linespread{0.94}\selectfont
Try every puzzle first. Where an answer shows a calculation, check the steps, not just the result.

\subhead{Chapter 1}
\textbf{Multix on the one way.} No. Each feature gives different groups and none is the right one. Which is best depends on what you want to find out. \textbf{Keep box.} Kept: the colour of each button. Set aside: size, shape and holes.

\textbf{1.1} $3+4+5=12$. Yes.\\
\textbf{1.2} $7+4=11$. One button was left out or not counted.\\
\textbf{1.3} The sort by colour made more groups, 3 against 2.\\
\textbf{1.4} Examples: size, colour, laces, material.\\
\textbf{1.5} Different features make different groups, and which sort is best depends on the question.\\
\textbf{1.6} The empty circle is 12.\\
\textbf{1.7} 4 have wings (A, B, E, F) and 4 do not.\\
\textbf{1.8} Yes. For example, the only gold button in the pile.\\
\textbf{1.9} In one sort a button is in one group only. It can be in ``red'' in a colour sort and ``four holes'' in a holes sort.\\
\textbf{1.10} Answers vary. Check the counts add up.

\subhead{Chapter 2}
\textbf{Multix on six legs.} No. A feature every bug has splits nothing. \textbf{Keep box.} Kept: whether a bug has wings. Set aside: spots and antennae.

\textbf{2.1} A, B, E, F.\\
\textbf{2.2} No. All eight bugs answer yes, so they stay in one group.\\
\textbf{2.3} A, B, G, H.\\
\textbf{2.4} It is an opinion, and people disagree.\\
\textbf{2.5} ``Is it longer than 2~cm?''\\
\textbf{2.6} 4: A, C, E, G.\\
\textbf{2.7} A and E.\\
\textbf{2.8} Spots do not change whether a bug has wings.\\
\textbf{2.9} Not always. A tiny or folded wing may be unclear, so say what counts as a wing.\\
\textbf{2.10} Answers vary.

\subhead{Chapter 3}
\textbf{Multix on three questions.} No. Eight bugs may not include one with an unexpected feature. \textbf{Keep box.} Kept: the path of answers. Set aside: other differences.

\textbf{3.1} C.\\
\textbf{3.2} 4. 8.\\
\textbf{3.3} 4.\\
\textbf{3.4} Example: ``Does it have a tail?''\\
\textbf{3.5} Every bug answers yes, so it does not split the group.\\
\textbf{3.6} Yes. Every combination is still covered.\\
\textbf{3.7} 4 and 4.\\
\textbf{3.8} 2.\\
\textbf{3.9} Answers vary.\\
\textbf{3.10} One.

\subhead{Chapter 4}
\textbf{Multix on trading.} No. The maker knows what each question means. Someone else must try it. \textbf{Keep box.} Kept: the idea of size. Set aside: the fuzzy line.

\textbf{4.1} 2.\\
\textbf{4.2} 2 more.\\
\textbf{4.3} ``Is it big?''\\
\textbf{4.4} She told her friend the answers, so the friend was not using the key on her own.\\
\textbf{4.5} $7+9+8=24$ and $24\div3=8$.\\
\textbf{4.6} Answers vary. Check the old question is circled and a clear one is outside.\\
\textbf{4.7} Neither is necessarily wrong. Compare how clear the questions are and how many people agree.\\
\textbf{4.8} (a) C (b) F (c) C (d) F.\\
\textbf{4.9} If they know the answers, their sorting may only copy them, which does not test the key.\\
\textbf{4.10} Answers vary.

\subhead{Chapter 5}
\textbf{Multix on the dog.} The key is not broken. The ``cat'' end really means fur without feathers, and the key only knew a cat. \textbf{Keep box.} Kept: groups that nest. Set aside: details of each animal.

\textbf{5.1} 1 and 2.\\
\textbf{5.2} 5.\\
\textbf{5.3} 1, 3, 4, 6 and 7.\\
\textbf{5.4} Butterfly and ant, 2 questions each.\\
\textbf{5.5} 2 have six legs and 6 do not.\\
\textbf{5.6} It ends at cat. The ``cat'' group is really fur without feathers.\\
\textbf{5.7} No. A spider has eight legs and the key puts it in the eight-legs group.\\
\textbf{5.8} Answers vary, for example goldfish in water, sparrow in air, snail on land. These groups do not match the key's.\\
\textbf{5.9} Answers vary.\\
\textbf{5.10} 2 and 6.

\subhead{Chapter 6}
\textbf{Multix on the worm.} No. The key had no question for a worm, so it was forced into the frog group. Change the key. \textbf{Keep box.} Kept: the old key, wrapped. Set aside: the wrong question.

\textbf{6.1} Bird. A bat has fur and feeds its young milk.\\
\textbf{6.2} ``Does it have legs?''\\
\textbf{6.3} Change the key. Forcing the fit hides the problem.\\
\textbf{6.4} No. ``Neither long nor short'' is unclear. Add a rule such as ``longer than 1~cm''.\\
\textbf{6.5} Answers vary. Check the old key is wrapped by the new question.\\
\textbf{6.6} 8.\\
\textbf{6.7} H.\\
\textbf{6.8} (a) S (b) G (c) B.\\
\textbf{6.9} Feathers.\\
\textbf{6.10} 5, since $2^5=32$.\\
\textbf{6.11} (a) M (b) R (c) M.\\
\textbf{6.12} No. Only ten things were tried, and new things may break it.\\
\textbf{6.13} Ask her to trade keys with a friend and compare.\\
\textbf{6.14 to 6.16} Answers vary. Check each split matches its labels.\\
\textbf{Slop check.} (a) You cannot tell: no title, no labels. (b) Many colours, stars, joke text. (c) A title, the colours and the counts with units. (d) The stars and jokes.

\clearpage\label{notes}
{\sans\bfseries\LARGE\color{ink} Words}\par
\small
\begin{multicols}{2}
\word{feature}{one thing you can check about an item, such as wings or colour.}
\word{group}{things that share a feature.}
\word{key}{a list of yes-or-no questions that sorts something into a group.}
\word{model}{a simple stand-in that keeps some things and drops others.}
\word{nested}{one group inside another.}
\word{specimen}{one example of a thing you are sorting.}
\end{multicols}

\vspace{0.4em}
{\sans\bfseries\large\color{sky!70!black} Notes for grown-ups}\par
This book teaches that a classification is a model: it keeps the features chosen for sorting and drops the rest. A key is tested by whether a second person gets the same answers, and a key is changed, not forced, when a surprise specimen does not fit. Each chapter has one activity, one worked example with made-up practice numbers, one trust box, one ``what did we keep?'' prompt and a page of puzzles.

\textbf{Materials and safety.} Buttons or small objects, leaves or shells, paper, a ruler, a friend and a safe place to look at animals. Do not touch or pick up animals, and do not touch plants you cannot name. Keep small objects away from young children and pets.

\textbf{The eight bugs} are made up so that every combination of three features appears once. They are not real insects. The animal key in Chapter~5 is a teaching key, and real classification uses many more features.

\textbf{The practice numbers} were made up so the arithmetic is easy to check.

\subhead{Sources}
\begin{itemize}[leftmargin=1.3em,itemsep=0.2em]
\item N. A. Campbell, J. B. Reece and others. \emph{Campbell Biology}. Pearson. For classification and the grouping of living things.
\item J. R. Taylor. \emph{An Introduction to Error Analysis}, 2nd edition. University Science Books, 1997. For agreement between independent readings.
\end{itemize}

\clearpage
\thispagestyle{blank}
\begin{tikzpicture}[remember picture,overlay]
  \fill[sun!35] ($(current page.south west)+(0.4in,0.4in)$) rectangle ($(current page.north east)-(0.4in,0.4in)$);
  \node[anchor=north,align=center,font=\sffamily\bfseries,text=ink,text width=5.4in] at ($(current page.north)+(0,-1.4in)$) {\Large Next in Science\\[6pt]\fontsize{28}{30}\selectfont Where the Energy Goes};
  \node[align=center,font=\sffamily,text width=4.6in,text=ink] at (current page.center) {A ball rolls down a ramp, across the floor, and stops. Book 8 follows the ball's energy from start to finish and keeps a running total of where it went.};
  \node at ($(current page.south)+(0,2.6in)$) {\scalebox{3}{\multixbot[ok]}};
  \node[font=\sffamily\bfseries,text=ink] at ($(current page.south)+(0,0.9in)$) {Ages 9 to 11 \quad$\star$\quad Science Book 7};
\end{tikzpicture}
